\section{Du serveur multimédia au homelab d'expérimentation}

Ce projet occupe une place particulière dans ce dossier. Il ne s'agit ni d'une mission réalisée au SHOM, ni d'un projet scolaire imposé. Il correspond à un environnement personnel que je fais évoluer depuis plusieurs années et que j'utilise comme terrain d'expérimentation pour les technologies de déploiement, d'orchestration et de supervision découvertes en formation. Son intérêt dans le dossier est donc complémentaire : il ne remplace pas les projets professionnels, mais me permet d'illustrer concrètement des pratiques modernes de \terme{devops}, de \terme{ci}, de \terme{cd} et de \terme{gitops}.

Le serveur a d'abord répondu à un besoin très simple : héberger Jellyfin pour disposer d'un serveur multimédia personnel. J'y ai ensuite ajouté Nextcloud afin de synchroniser mon dossier de documents entre ma tour et mon ordinateur portable, tous deux sous Manjaro. Progressivement, la machine est devenue un environnement permanent de test et de déploiement. Lorsqu'une technologie abordée en cours pouvait être expérimentée sur une infrastructure réelle, je cherchais à l'intégrer ou à la tester sur ce serveur plutôt que de limiter l'exercice à une machine virtuelle jetable.

Cette démarche a notamment conduit à l'utilisation de Docker, puis de Kubernetes avec k3s, à la mise en place d'un runner GitHub Actions auto-hébergé, à différentes expérimentations de déploiement avec Terraform, Argo CD et Helm, puis à l'installation d'une stack d'observabilité composée de Prometheus, Grafana, Loki et Grafana Alloy. L'objectif n'était pas de construire une plateforme de production hautement disponible, mais de disposer d'un environnement suffisamment réaliste pour confronter les outils étudiés à des contraintes concrètes d'exploitation.

\section{Compétences principalement mobilisées}

Le projet renforce principalement les compétences du bloc 3 liées à l'intégration, au déploiement et aux opérations continues. Il apporte également un exemple de veille et d'appropriation technologique réalisée en dehors d'un cadre imposé.

\begin{table}[H]
\centering
\begin{tabular}{p{2.2cm} p{11.8cm}}
\toprule
\textbf{Compétence} & \textbf{Mobilisation dans le homelab} \\
\midrule
C1.2 & Expérimentation de technologies découvertes en formation et comparaison progressive de plusieurs modes de déploiement : Docker, Kubernetes, Terraform, Argo CD et Helm. \\
C3.1 & Utilisation de GitHub Actions et d'un runner auto-hébergé pour construire des images, publier des artefacts et déclencher les opérations de déploiement. \\
C3.3 & Mise sous gestion Git de la stack d'observabilité et synchronisation automatique de son état avec Argo CD. \\
C3.4 & Déploiement automatisé d'une application conteneurisée avec GitHub Actions, GHCR, Helm et k3s. \\
C3.5 & Mise en place d'une stack Prometheus/Grafana/Loki/Alloy pour observer les ressources Kubernetes, le serveur et l'activité HTTP du reverse proxy. \\
C3.6 & Formalisation de l'architecture, des procédures d'installation et des commandes de vérification afin de rendre l'environnement compréhensible et reproductible. \\
\bottomrule
\end{tabular}
\caption{Compétences mobilisées dans le projet de homelab}
\label{tab:competences-homelab}
\end{table}

\section{Une infrastructure personnelle contrainte mais représentative}

Le serveur repose sur un ancien ordinateur portable Asus VivoBook récupéré et réemployé plutôt que sur une machine achetée spécifiquement pour cet usage. Il est équipé d'un AMD Ryzen 7 3700U, d'environ 18~Go de mémoire vive, d'un SSD interne de 500~Go et d'un second SSD externe de 500~Go. Le système utilisé est Ubuntu Server 24.04 LTS, sans environnement graphique.

Ce matériel est suffisamment dimensionné pour héberger plusieurs services et un cluster k3s mono-nœud, mais il impose de rester proportionné dans les choix d'architecture. Le cluster n'est pas conçu pour assurer de la haute disponibilité et l'ensemble de l'infrastructure dépend physiquement d'une seule machine. En contrepartie, ce contexte rend les expérimentations proches de contraintes réelles : ressources limitées, stockage local, exposition réseau, certificats, supervision et coexistence de plusieurs modes d'exécution.

Tous les services ne sont volontairement pas placés dans Kubernetes. Jellyfin fonctionne directement sur l'hôte afin de simplifier l'utilisation de l'accélération matérielle. Nextcloud et MariaDB sont regroupés dans un Docker Compose dédié. Caddy possède également son propre Docker Compose et joue le rôle de reverse proxy central. Le cluster k3s est utilisé pour les expérimentations Kubernetes, la stack d'observabilité et certaines applications, dont PasswordGenerator. Enfin, le runner GitHub Actions fonctionne directement sur le serveur afin d'exécuter les commandes nécessitant un accès local au cluster.

\begin{figure}[H]
\centering
\fbox{\parbox{0.90\textwidth}{\centering
Emplacement prévu : schéma global du homelab.\\[0.15cm]
Internet $\rightarrow$ Cloudflare $\rightarrow$ Caddy $\rightarrow$ services locaux.\\
Caddy distribue vers Jellyfin en bare metal, Nextcloud/MariaDB sous Docker Compose et le cluster k3s.\\
Dans k3s : Argo CD, PasswordGenerator et le namespace observability avec Prometheus, Grafana, Loki et Alloy.\\
Le runner GitHub Actions est installé directement sur le serveur.}}
\caption{Architecture générale du homelab personnel}
\label{fig:homelab-architecture}
\end{figure}

Cette architecture hybride n'est pas le résultat d'un plan initial figé. Elle traduit l'évolution progressive du serveur. Les services historiques ont été conservés dans le mode d'exécution qui répondait correctement au besoin, tandis que Kubernetes a été introduit pour expérimenter l'orchestration et les mécanismes de déploiement. Cela m'a permis de comparer les outils sans chercher à imposer Kubernetes à tous les usages.

\section{Exposer les services malgré le CGNAT}

L'accès depuis Internet introduit une contrainte supplémentaire : ma connexion est située derrière un CGNAT et ne dispose donc pas d'une adresse IPv4 publique directement joignable. Le serveur possède en revanche une adresse IPv6 stable. Le domaine \texttt{thomasgraber.fr} est géré chez Cloudflare ; les sous-domaines utilisés par les services possèdent un enregistrement AAAA vers l'adresse IPv6 du serveur avec le proxy Cloudflare activé.

Le trafic HTTP et HTTPS passe ainsi par Cloudflare avant d'atteindre le serveur en IPv6. Le mode SSL/TLS utilisé est \textit{Full}, ce qui conserve une connexion chiffrée entre le client et Cloudflare, puis entre Cloudflare et le serveur d'origine. Caddy reçoit les requêtes et les relaie vers les services locaux en fonction du nom de domaine demandé.

Par exemple, \texttt{password.thomasgraber.fr} est redirigé vers le NodePort Kubernetes du générateur de mots de passe, \texttt{nc.thomasgraber.fr} vers Nextcloud et \texttt{jellyfin.thomasgraber.fr} vers le service Jellyfin exécuté sur l'hôte. Les interfaces Grafana, Prometheus et Argo CD peuvent être exposées selon le même principe. Cette configuration m'a permis de manipuler concrètement des sujets qui dépassent le seul développement applicatif : résolution DNS, IPv6, reverse proxy, TLS et routage vers plusieurs modes d'hébergement.

\section{PasswordGenerator comme support d'expérimentation du déploiement}

PasswordGenerator est une petite application web volontairement simple. Sa logique fonctionnelle n'est pas l'élément intéressant dans ce chapitre. Elle m'a surtout servi de support stable pendant plusieurs cours afin de tester différentes stratégies de livraison sans devoir reconstruire une application métier à chaque expérimentation.

L'historique du dépôt permet de retrouver une évolution significative des modes de déploiement. L'application a d'abord été conteneurisée avec Docker puis exécutée sur Kubernetes. J'ai ensuite expérimenté Terraform pour déclarer le namespace, le Deployment et le Service Kubernetes depuis une configuration d'infrastructure. Une étape suivante a consisté à remplacer cette gestion directe par des manifestes versionnés et Argo CD, afin de tester une approche GitOps. Enfin, le déploiement actuellement utilisé repose sur un chart Helm déclenché depuis GitHub Actions.

Ces étapes ne correspondent pas à cinq architectures totalement distinctes. Docker reste utilisé pour construire le conteneur et Kubernetes reste la plateforme d'exécution. Ce qui évolue surtout est la manière de décrire et d'appliquer l'état souhaité. Cette succession m'a permis de comparer plusieurs niveaux d'abstraction : description d'infrastructure avec Terraform, ressources Kubernetes explicites sous GitOps, puis packaging avec Helm.

\subsection{Pipeline actuellement utilisé}

Le workflow actuel se déclenche lors d'un push sur la branche principale. Il s'exécute sur le runner auto-hébergé, se connecte à GitHub Container Registry, construit une image Docker et la publie avec le SHA du commit comme tag. Cette stratégie permet d'associer précisément une version déployée à un état du dépôt. Le workflow exécute ensuite un \texttt{helm upgrade --install} en fournissant au chart le repository de l'image et le SHA courant.

\begin{figure}[H]
\centering
\fbox{\parbox{0.90\textwidth}{\centering
Emplacement prévu : chaîne CI/CD PasswordGenerator.\\[0.15cm]
Push sur \texttt{main} $\rightarrow$ GitHub Actions auto-hébergé $\rightarrow$ build Docker $\rightarrow$ GHCR\\
$\rightarrow$ \texttt{helm upgrade --install} avec le SHA du commit $\rightarrow$ k3s $\rightarrow$ NodePort $\rightarrow$ Caddy.}}
\caption{Pipeline de déploiement actuel de PasswordGenerator}
\label{fig:password-cicd}
\end{figure}

Le chart configure notamment plusieurs réplicas, un service NodePort ainsi que des probes de disponibilité et de vivacité. L'application est ensuite rendue accessible publiquement via Caddy et Cloudflare. Le pipeline automatise donc effectivement la construction, la publication et le déploiement. Il comporte cependant une limite importante : l'application étant très simple et ayant principalement servi de support d'expérimentation, aucune suite de tests automatisés n'est actuellement exécutée avant la mise en production. Cet élément distingue une chaîne technique fonctionnelle d'un pipeline pleinement industrialisé.

\section{Mettre en place l'observabilité du cluster}

La stack d'observabilité est issue d'un enseignement consacré à Prometheus, Grafana et aux outils associés. Après avoir découvert leur rôle en cours, j'ai choisi de ne pas conserver l'exercice dans un environnement temporaire et de l'intégrer à mon propre serveur. La démarche était donc d'abord pédagogique, mais l'environnement obtenu fonctionne aujourd'hui en continu sur le homelab.

J'ai commencé par installer manuellement les composants avec Helm afin de comprendre les charts utilisés et de valider les fichiers de valeurs. Prometheus et Grafana sont fournis par \texttt{kube-prometheus-stack}, qui installe également Alertmanager, kube-state-metrics et node-exporter. Loki est utilisé pour le stockage et l'interrogation des logs. Grafana Alloy assure la collecte de logs et d'événements Kubernetes et les transmet à Loki.

Les choix sont adaptés au contexte mono-nœud. Loki utilise un déploiement monolithique avec une seule réplique et un stockage persistant local. Prometheus conserve les métriques pendant sept jours avec une limite de taille de 6~Go. Ces réglages évitent de dimensionner l'installation comme une plateforme distribuée alors que le besoin est celui d'un homelab personnel.

Grafana fournit les tableaux de bord Kubernetes installés avec la stack : ressources des pods et workloads, kubelet, réseau par cluster ou namespace, volumes persistants, métriques des nœuds et vue d'ensemble de Prometheus. Dans mon usage actuel, je consulte principalement la consommation du serveur, l'état des workloads et l'activité réseau. La stack ne constitue pas encore un dispositif de supervision proactif complet : Alertmanager est installé mais aucune chaîne de notification n'est configurée.

\subsection{Superviser également Caddy}

Prometheus n'observe pas uniquement Kubernetes. J'ai activé l'exposition des métriques dans Caddy et ajouté un scrape spécifique dans la configuration Prometheus. Prometheus interroge ainsi périodiquement l'endpoint \texttt{/metrics} de Caddy et peut intégrer son activité HTTP dans les données consultées depuis Grafana.

Alloy complète cette approche pour les logs. Il récupère notamment des fichiers système de l'hôte et découvre les conteneurs Docker à partir du socket Docker monté en lecture. Les logs des conteneurs peuvent ensuite être enrichis par des labels avant leur envoi vers Loki. Pour Caddy, un traitement spécifique extrait plusieurs champs structurés, par exemple le statut HTTP, l'hôte demandé et le niveau du message.

Cette partie a demandé plusieurs ajustements. Une première configuration basée directement sur les fichiers de logs Docker ne faisait pas apparaître correctement les conteneurs. Je l'ai remplacée par la découverte Docker d'Alloy via le socket de l'hôte et par des règles de relabellisation. Cette correction est représentative de l'intérêt du serveur personnel : au-delà de la configuration théorique d'un outil, il faut comprendre pourquoi les données attendues n'arrivent pas et adapter la collecte au contexte réel.

\section{Passer de Helm manuel à une gestion GitOps}

Après avoir validé manuellement la stack d'observabilité, j'ai placé les fichiers \texttt{values} et les objets \texttt{Application} Argo CD dans un dépôt GitHub dédié. Le but était que Git devienne la référence décrivant l'état attendu de la supervision.

Les objets Argo CD utilisent deux sources. La première est le repository Helm officiel correspondant au composant à déployer. La seconde est mon dépôt Git, dans lequel sont stockés les fichiers de configuration spécifiques au homelab. Argo CD récupère ainsi le chart officiel mais utilise les valeurs versionnées dans mon dépôt.

Les politiques de synchronisation sont automatisées avec \texttt{prune} et \texttt{selfHeal}. Une ressource supprimée de la déclaration peut donc être nettoyée, tandis qu'un écart entre l'état du cluster et l'état attendu peut être corrigé par Argo CD. Comme le dépôt est privé, l'accès d'Argo CD repose sur une clé SSH dédiée, configurée côté GitHub comme deploy key et côté Kubernetes sous forme de Secret.

Un workflow GitHub Actions complète ce dispositif. Il ne déploie pas directement Prometheus, Grafana ou Loki. Lorsqu'un changement concerne les fichiers Argo CD ou les fichiers de valeurs, le runner applique les objets \texttt{Application}, force leur rafraîchissement puis affiche leur état. Argo CD reste ensuite responsable de la synchronisation avec Kubernetes.

\begin{figure}[H]
\centering
\fbox{\parbox{0.90\textwidth}{\centering
Emplacement prévu : comparaison des deux modèles utilisés.\\[0.15cm]
\textbf{PasswordGenerator :} GitHub Actions pousse directement le nouvel état via Helm.\\
\textbf{Observabilité :} Git décrit l'état désiré, GitHub Actions actualise les Applications Argo CD, puis Argo CD réconcilie le cluster.}}
\caption{Deux stratégies de déploiement expérimentées sur le même homelab}
\label{fig:homelab-cd-gitops}
\end{figure}

La coexistence de ces deux modèles constitue l'un des principaux apports de ce projet. Dans le premier cas, la CI/CD exécute directement l'opération de déploiement. Dans le second, le pipeline met à disposition une déclaration et laisse un contrôleur présent dans le cluster comparer en continu l'état réel avec l'état attendu.

\section{Valider le fonctionnement GitOps par une modification observable}

Pour vérifier que le mécanisme ne fonctionnait pas uniquement en théorie, j'ai volontairement modifié dans Git le NodePort utilisé par Grafana, de 32431 à 32432. Le push a déclenché GitHub Actions. Les Applications Argo CD ont été rafraîchies puis ont retrouvé un état \texttt{Synced} et \texttt{Healthy}. Les vérifications Kubernetes ont confirmé que Grafana utilisait désormais le nouveau port et qu'aucun deuxième Service ou Deployment n'avait été créé.

Ce test a également mis en évidence un effet de bord intéressant : Caddy pointait toujours vers l'ancien NodePort. Le service restait temporairement accessible tant que la configuration active de Caddy n'avait pas été relue, puis devenait inaccessible après redémarrage du reverse proxy. Il a fallu mettre à jour le Caddyfile pour rétablir la cohérence entre l'exposition externe et le service Kubernetes.

Cette situation montre une limite du périmètre GitOps actuel : Argo CD garantit l'état des ressources qu'il gère dans Kubernetes, mais Caddy est exploité séparément avec Docker Compose. Une modification d'interface entre les deux environnements peut donc nécessiter une adaptation manuelle. L'automatisation d'un sous-système ne supprime pas les dépendances avec les composants placés hors de son périmètre.

\section{Difficultés rencontrées}

Plusieurs difficultés rencontrées pendant la mise en place ont participé à l'apprentissage. La configuration d'un runner auto-hébergé a d'abord posé des problèmes de droits sur son répertoire d'installation. Il a fallu distinguer les opérations nécessitant les privilèges système de la configuration du runner lui-même et corriger le propriétaire du dossier.

Argo CD a également produit une erreur de comparaison tant qu'il ne disposait pas des credentials nécessaires pour lire le dépôt privé contenant les fichiers de valeurs. La mise en place d'une deploy key SSH dédiée a permis de conserver le dépôt privé sans placer de clé dans Git.

Enfin, l'installation de \texttt{kube-prometheus-stack} sur k3s faisait apparaître comme indisponibles certains endpoints que cette distribution n'expose pas de la même manière qu'un cluster Kubernetes classique. J'ai désactivé les cibles concernées afin de ne pas conserver des indicateurs \texttt{DOWN} artificiels. Ces adaptations rappellent qu'un chart Helm fournit une base générique mais ne dispense pas de comprendre l'environnement dans lequel il est installé.

\section{Limites et prise de recul}

Le homelab reste avant tout un environnement personnel d'apprentissage. Plusieurs choix seraient insuffisants pour une infrastructure de production critique. Le cluster k3s ne comporte qu'un seul nœud ; la multiplication des réplicas d'une application améliore certains scénarios de redémarrage mais ne protège pas contre la panne physique de la machine. Le stockage \texttt{local-path} lie également les volumes au nœud unique.

Le runner GitHub Actions est installé sur le même serveur que l'environnement qu'il déploie. Cette organisation est pratique dans un homelab mais crée un couplage fort entre l'outil de livraison et la cible. De même, certaines interfaces techniques sont exposées derrière le reverse proxy et nécessiteraient une politique d'accès plus restrictive dans un contexte professionnel. La configuration Caddy utilisée pour joindre Argo CD contient par ailleurs une désactivation de la vérification TLS du certificat de l'upstream local, simplification qu'il serait préférable de supprimer dans une architecture plus stricte.

L'observabilité reste majoritairement passive. Je consulte les métriques et l'activité réseau, mais Alertmanager n'est pas encore configuré pour notifier une anomalie. Une évolution logique serait de définir quelques alertes réellement utiles, par exemple sur l'espace disque, la mémoire disponible, l'indisponibilité d'un service ou les erreurs HTTP, puis de configurer un canal de notification.

La stratégie de sauvegarde constitue une autre limite. Les documents Nextcloud sont synchronisés entre le serveur, ma tour et mon ordinateur portable. Cette réplication protège contre la perte isolée d'un équipement, mais elle ne remplace pas une sauvegarde indépendante et versionnée : une suppression ou une corruption synchronisée peut se propager aux différentes copies. Les composants Kubernetes et les bases de données ne disposent pas non plus aujourd'hui d'une politique de sauvegarde complète.

Enfin, le projet met en évidence une différence importante entre expérimentation et industrialisation. PasswordGenerator dispose d'une chaîne de déploiement automatisée, mais pas encore d'un plan de tests automatisés avant déploiement. La stack d'observabilité est gérée en GitOps, mais plusieurs services du homelab restent administrés séparément. Ces écarts ne remettent pas en cause l'intérêt du projet ; ils permettent au contraire d'identifier ce qui serait nécessaire pour passer d'un laboratoire personnel à une infrastructure professionnelle plus robuste.

\section{Bilan}

Ce homelab m'a permis d'approfondir des sujets que je manipule moins directement dans mes missions professionnelles : conteneurisation, orchestration Kubernetes, runners auto-hébergés, registres d'images, Helm, GitOps, collecte de métriques et centralisation des logs. La réutilisation du même environnement dans plusieurs enseignements a également rendu visibles les différences entre outils qui peuvent sembler similaires lorsqu'ils sont étudiés séparément.

L'évolution de PasswordGenerator illustre particulièrement cette démarche. Une application simple a servi de support à plusieurs modèles de déploiement successifs, tandis que la stack d'observabilité m'a permis de mettre en pratique une approche GitOps plus complète avec Argo CD. Le serveur est ainsi devenu moins un produit à terminer qu'un environnement de veille et d'expérimentation continue.

Le projet présente enfin un intérêt en matière de numérique responsable à une échelle modeste mais concrète. L'infrastructure repose sur le réemploi d'un ancien ordinateur portable et les composants sont dimensionnés pour un besoin personnel, par exemple avec des durées de rétention limitées et une architecture Loki monolithique. Je ne considère pas cela comme une démarche complète d'éco-conception, mais comme un choix cohérent avec le niveau de service recherché et les ressources réellement nécessaires.
